\documentclass[12pt,a4paper]{article}

\usepackage[brazilian]{babel}
\usepackage[T1]{fontenc}
\usepackage{amsmath,amssymb}
\usepackage{geometry}
\usepackage{booktabs}
\geometry{margin=2.5cm}

\title{Formulario de Formulas\\Capitulos 2, 3, 4 e 5}
\author{}
\date{\today}

\begin{document}
\maketitle

\section*{1. Base da Analise Dimensional}

\textbf{Dimensao geral:}
\begin{equation}
[Q] = M^a L^b T^c \Theta^d \cdots
\end{equation}

\textbf{Homogeneidade dimensional:}
\begin{equation}
A = B + C \;\Rightarrow\; [A]=[B]=[C]
\end{equation}

\textbf{Forma de potencia (ansatz):}
\begin{equation}
Y = C\,x_1^{\alpha_1}x_2^{\alpha_2}\cdots x_n^{\alpha_n}
\end{equation}

\textbf{Numero de grupos adimensionais (Buckingham Pi):}
\begin{equation}
N_\Pi = n - m
\end{equation}
onde $n$ = numero de variaveis e $m$ = numero de dimensoes fundamentais independentes.

\textbf{Forma final do teorema Pi:}
\begin{equation}
\Phi(\Pi_1,\Pi_2,\ldots,\Pi_{n-m})=0
\end{equation}

\section*{2. Dimensoes Comuns}
\begin{center}
\begin{tabular}{ll}
\toprule
Quantidade & Dimensao \\
\midrule
Velocidade $v$ & $LT^{-1}$ \\
Aceleracao $a$ & $LT^{-2}$ \\
Forca $F$ & $MLT^{-2}$ \\
Energia $E$ & $ML^2T^{-2}$ \\
Potencia $P$ & $ML^2T^{-3}$ \\
Pressao $p$ & $ML^{-1}T^{-2}$ \\
Densidade $\rho$ & $ML^{-3}$ \\
Viscosidade dinamica $\mu$ & $ML^{-1}T^{-1}$ \\
Frequencia $f$ & $T^{-1}$ \\
\bottomrule
\end{tabular}
\end{center}

\section*{3. Formulas Classicas do Capitulo}

\textbf{Pendulo simples (pequenas amplitudes):}
\begin{equation}
T_0 = 2\pi\sqrt{\frac{l}{g}}
\end{equation}

\textbf{Queda livre (escala dimensional):}
\begin{equation}
V = C\sqrt{gh}
\end{equation}
(com fisica completa: $V=\sqrt{2gh}$)

\textbf{Velocidade do som (escala):}
\begin{equation}
c = C\sqrt{\frac{p}{\rho}}
\end{equation}
(com gas ideal: $c=\sqrt{\gamma\,p/\rho}$)

\textbf{Arrasto viscoso (Stokes):}
\begin{equation}
F_D = 3\pi\mu V d
\end{equation}

\section*{4. Exemplo Pi -- Grua no Estadio}

\textbf{Grupos adimensionais:}
\begin{align}
\Pi_1 &= \frac{L}{h}, \\
\Pi_2 &= \frac{g}{h\omega^2}, \\
\Pi_3 &= \frac{\tau}{m h\omega^2}
\end{align}

\textbf{Escalas para a frequencia:}
\begin{align}
\omega &\propto \sqrt{\frac{g}{h}}, \\
\omega &\propto \sqrt{\frac{\tau}{m h}}
\end{align}

\section*{5. Principais Numeros Adimensionais}
\begin{align}
Re &= \frac{\rho U L}{\mu} & \text{(Reynolds)} \\
Fr &= \frac{U}{\sqrt{gL}} & \text{(Froude)} \\
Ma &= \frac{U}{c} & \text{(Mach)} \\
We &= \frac{\rho U^2 L}{\sigma} & \text{(Weber)} \\
Pe &= \frac{UL}{D} & \text{(Peclet)} \\
Nu &= \frac{hL}{k} & \text{(Nusselt)} \\
Pr &= \frac{\mu c_p}{k} & \text{(Prandtl)} \\
Bi &= \frac{hL}{k_s} & \text{(Biot)} \\
St &= \frac{fL}{U} & \text{(Stokes nao estacionario)} \\
Eu &= \frac{\Delta p}{\rho U^2} & \text{(Euler)}
\end{align}

\section*{6. Conversao de Unidades (Exemplo)}
\begin{equation}
65\,\frac{\text{mi}}{\text{hr}}
\cdot \frac{5280\,\text{ft}}{\text{mi}}
\cdot \frac{0.3048\,\text{m}}{\text{ft}}
\cdot \frac{1\,\text{hr}}{3600\,\text{s}}
= 29.06\,\text{m/s}
\end{equation}

\section*{7. Formulas do Capitulo 3 (Escala, Escalabilidade e Fractais)}

\textbf{Nao dimensionalizacao (transporte 1D):}
\begin{equation}
\frac{\partial c}{\partial t} + u\frac{\partial c}{\partial x} = D\frac{\partial^2 c}{\partial x^2}
\end{equation}
\begin{equation}
\hat{x}=\frac{x}{L},\quad \hat{t}=\frac{t}{T},\quad \hat{c}=\frac{c}{C_0},\quad \hat{u}=\frac{u}{U}
\end{equation}
com $T=L/U$:
\begin{equation}
\frac{\partial \hat{c}}{\partial \hat{t}} + \hat{u}\frac{\partial \hat{c}}{\partial \hat{x}}
= \frac{1}{Pe}\frac{\partial^2 \hat{c}}{\partial \hat{x}^2},
\qquad
Pe=\frac{UL}{D}
\end{equation}

\textbf{Lei de potencia:}
\begin{equation}
Y = kX^{\alpha}
\end{equation}
\begin{equation}
\log Y = \log k + \alpha\log X
\end{equation}

\textbf{Escalas geometricas:}
\begin{equation}
A \propto L^2,\qquad V \propto L^3
\end{equation}
Se $L' = \lambda L$:
\begin{equation}
A' = \lambda^2 A,\qquad V' = \lambda^3 V
\end{equation}

\textbf{Pendulo (escala):}
\begin{equation}
T = 2\pi\sqrt{\frac{\ell}{g}},
\qquad
T\propto \ell^{1/2}
\end{equation}

\textbf{Escalabilidade computacional:}
\begin{equation}
C(N) = aN^2,
\qquad
C(2N)=4C(N)
\end{equation}

\textbf{Escala termica:}
\begin{equation}
\dot{Q}\sim hA\Delta T,\qquad \dot{Q}\propto L^2
\end{equation}
\begin{equation}
E\propto L^3,
\qquad
\tau\sim\frac{E}{\dot{Q}}\propto L
\end{equation}

\textbf{Lei de Kleiber (alometria):}
\begin{equation}
q_0 \approx 70\,M^{3/4}
\end{equation}

\textbf{Dimensao fractal por auto-similaridade:}
\begin{equation}
D = \frac{\log N}{\log(1/r)}
\end{equation}

\textbf{Fractais classicos:}
\begin{equation}
D_{\text{Koch}} = \frac{\log 4}{\log 3} \approx 1.2619,
\qquad
D_{\text{Sierpinski}} = \frac{\log 3}{\log 2} \approx 1.585
\end{equation}

\textbf{Forma usada na prova (fractal auto-similar):}
\begin{equation}
D_f = \frac{\ln N_p}{\ln(1/p)}
\end{equation}

\section*{8. Formulas do Capitulo 4 (Expansao de Taylor)}

\subsection*{Definicao e forma geral}

\textbf{Serie de Taylor em torno de $a$:}
\begin{equation}
f(x) = \sum_{k=0}^{n} \frac{f^{(k)}(a)}{k!}(x-a)^k + R_{n+1}(x)
\end{equation}

\textbf{Serie de Maclaurin (caso $a=0$):}
\begin{equation}
f(x) = \sum_{k=0}^{n} \frac{f^{(k)}(0)}{k!}x^k + R_{n+1}(x)
\end{equation}

\textbf{Resto de Lagrange:}
\begin{equation}
R_{n+1}(x) = \frac{f^{(n+1)}(\xi)}{(n+1)!}(x-a)^{n+1}, \quad \xi \text{ entre } a \text{ e } x
\end{equation}

\subsection*{Series fundamentais (Maclaurin)}

\textbf{Exponencial:}
\begin{equation}
e^x = 1 + x + \frac{x^2}{2!} + \frac{x^3}{3!} + \cdots
\end{equation}

\textbf{Seno e cosseno:}
\begin{align}
\sin x &= x - \frac{x^3}{3!} + \frac{x^5}{5!} - \cdots \\
\cos x &= 1 - \frac{x^2}{2!} + \frac{x^4}{4!} - \cdots
\end{align}

\textbf{Logaritmo natural (valido para $|x|<1$):}
\begin{equation}
\ln(1+x) = x - \frac{x^2}{2} + \frac{x^3}{3} - \frac{x^4}{4} + \cdots
\end{equation}

\textbf{Serie binomial (valida para $|x|<1$, $n \in \mathbb{R}$):}
\begin{equation}
(1+x)^n = 1 + nx + \frac{n(n-1)}{2!}x^2 + \frac{n(n-1)(n-2)}{3!}x^3 + \cdots
\end{equation}

\subsection*{Aproximacoes praticas}

\textbf{Aproximacao binomial (primeira ordem):}
\begin{equation}
(1+x)^n \approx 1 + nx, \quad |x| \ll 1
\end{equation}

Casos frequentes:
\begin{align}
\sqrt{1+x} &\approx 1 + \frac{x}{2} \\
(1+x)^{-1} &\approx 1 - x \\
(1+x)^{-3/2} &\approx 1 - \frac{3}{2}x
\end{align}

\textbf{Angulos pequenos:}
\begin{equation}
\sin\theta \approx \theta, \quad \cos\theta \approx 1 - \frac{\theta^2}{2}, \quad |\theta| \ll 1
\end{equation}

\subsection*{Linearizacao de modelos nao lineares}

\textbf{EDO nao linear perto de equilibrio $x^*$ (com $f(x^*)=0$):}
\begin{equation}
\dot{x} = f(x) \approx f'(x^*)(x-x^*), \quad \eta = x - x^*
\end{equation}
\begin{equation}
\dot{\eta} \approx f'(x^*)\eta \quad \text{(modelo linearizado)}
\end{equation}

\subsection*{Aplicacoes em modelagem}

\textbf{Pendulo simples:}
\begin{equation}
\ddot{\theta} + \frac{g}{\ell}\sin\theta = 0 \quad \Rightarrow \quad
\ddot{\theta} + \frac{g}{\ell}\theta = 0, \quad |\theta| \ll 1
\end{equation}
\begin{equation}
T = 2\pi\sqrt{\frac{\ell}{g}}
\end{equation}

\textbf{Modelo logistico:}
\begin{equation}
\dot{P} = rP\left(1 - \frac{P}{K}\right)
\end{equation}
Equilibrios: $P^* = 0$ (instavel) e $P^* = K$ (estavel se $r>0$).

\textbf{Decaimento exponencial (ex: C-14):}
\begin{equation}
C(t) = C_0 e^{-\lambda t}, \quad \lambda = \frac{\ln 2}{t_{1/2}}
\end{equation}
\begin{equation}
\delta t = \left|\frac{dt}{dN}\right|\Delta N
= \left|\frac{dt}{dC}\right|\Delta C
= \frac{t_{1/2}}{\ln 2\cdot C}\,\Delta C
\end{equation}

\textbf{Depreciacacao exponencial:}
\begin{equation}
V(t) = V_0 (1-r)^t = V_0 e^{-\lambda t}, \quad \lambda = -\ln(1-r) \approx r, \quad r \ll 1
\end{equation}

\section*{9. Formulas do Capitulo 5 (Comportamento Exponencial)}

\subsection*{Modelo exponencial basico}

\textbf{Crescimento/decaimento proporcional ao estado:}
\begin{equation}
\frac{dy}{dt} = ky
\end{equation}
\begin{equation}
y(t) = y_0 e^{kt}
\end{equation}

\textbf{Forma de decaimento (com $\lambda>0$):}
\begin{equation}
\frac{dy}{dt} = -\lambda y
\quad\Rightarrow\quad
y(t) = y_0 e^{-\lambda t}
\end{equation}

\textbf{Constante de tempo e meia-vida:}
\begin{equation}
\tau = \frac{1}{\lambda},
\qquad
t_{1/2} = \frac{\ln 2}{\lambda},
\qquad
\lambda = \frac{\ln 2}{t_{1/2}}
\end{equation}

\subsection*{Linearizacao em escala logaritmica}

\textbf{Reta em grafico semilog:}
\begin{equation}
\ln y = \ln y_0 + kt
\end{equation}
Inclinacao da reta: $k$.

\subsection*{Modelos classicos do capitulo}

\textbf{Juros compostos continuos:}
\begin{equation}
\frac{dC}{dt} = rC
\quad\Rightarrow\quad
C(t) = C_0 e^{rt}
\end{equation}

\textbf{Decaimento radioativo:}
\begin{equation}
\frac{dN}{dt} = -\lambda N
\quad\Rightarrow\quad
N(t) = N_0 e^{-\lambda t}
\end{equation}

\textbf{Resfriamento de Newton:}
\begin{equation}
\frac{dT}{dt} = -k(T-T_a)
\quad\Rightarrow\quad
T(t)=T_a + (T_0-T_a)e^{-kt}
\end{equation}

\textbf{Circuito RC (descarga):}
\begin{equation}
\frac{dV}{dt} = -\frac{1}{RC}V
\quad\Rightarrow\quad
V(t)=V_0 e^{-t/(RC)},
\qquad \tau = RC
\end{equation}

\subsection*{Modelo logistico (Verhulst)}

\textbf{Equacao diferencial:}
\begin{equation}
\frac{dP}{dt} = rP\left(1-\frac{P}{K}\right)
\end{equation}

\textbf{Solucao com $P(0)=P_0$:}
\begin{equation}
P(t) = \frac{K}{1 + \left(\dfrac{K-P_0}{P_0}\right)e^{-rt}}
\end{equation}

\textbf{Ponto de inflexao da curva logistica:}
\begin{equation}
t^* = \frac{1}{r}\ln\left(\frac{K-P_0}{P_0}\right)
\end{equation}
e ocorre quando $P=K/2$.

\subsection*{Equacoes de Lanchester}

\textbf{Sistema acoplado:}
\begin{equation}
\frac{dE}{dt} = -\lambda_F F,
\qquad
\frac{dF}{dt} = -\lambda_E E
\end{equation}

\textbf{Invariante (lei quadratica):}
\begin{equation}
\lambda_E E^2 - \lambda_F F^2 = \lambda_E E_0^2 - \lambda_F F_0^2 = \text{constante}
\end{equation}

\subsection*{Propagacao de incerteza no decaimento}

Para $t = -\ln(C/C_0)/\lambda$:
\begin{equation}
\delta t = \left|\frac{dt}{dN}\right|\Delta N
\end{equation}
onde, para o caso de C-14 ($N\equiv C$):
\begin{equation}
\delta t = \left|\frac{dt}{dC}\right|\Delta C
= \frac{1}{\lambda C}\,\Delta C
= \frac{t_{1/2}}{\ln 2\cdot C}\,\Delta C
\end{equation}

\end{document}
